% !TEX root = Resume_for_Frehers.tex
%%%%%%%%%%%%%%%%%%%%%%%%%%%%%%%%%%%%%%%%%
% Medium Length Professional CV
% LaTeX Template
% Version 2.0 (8/5/13)
%
% This template has been downloaded from:
% http://www.LaTeXTemplates.com
%
% Original author:
% Rishi Shah 
%
% Important note:
% This template requires the resume.cls file to be in the same directory as the
% .tex file. The resume.cls file provides the resume style used for structuring the
% document.
%
%%%%%%%%%%%%%%%%%%%%%%%%%%%%%%%%%%%%%%%%%

%----------------------------------------------------------------------------------------
%	PACKAGES AND OTHER DOCUMENT CONFIGURATIONS
%----------------------------------------------------------------------------------------

\documentclass{resume} % Use the custom resume.cls style

\usepackage{luatexja}
\usepackage{bibentry}
\usepackage[left=0.75in,top=0.6in,right=0.75in,bottom=0.6in]{geometry} % Document margins
\newcommand{\tab}[1]{\hspace{.2667\textwidth}\rlap{#1}}
\newcommand{\itab}[1]{\hspace{0em}\rlap{#1}}
\name{Arata Horie} % Your name
\address{507 (5F), CCR Bldg., 4-6-1 KOMABA MEGURO-KU, TOKYO 153-8505, JAPAN} % Your address
%\address{123 Pleasant Lane \\ City, State 12345} % Your secondary addess (optional)
\address{ \\ a.horie@commissure.co.jp} % Your phone number and email

\begin{document}
\nobibliography{horie}
\bibliographystyle{unsrt}

%----------------------------------------------------------------------------------------
%	EDUCATION SECTION
%----------------------------------------------------------------------------------------


\begin{rSection}{Education}

{\bf The University of Tokyo} \hfill {\bf Tokyo, Japan} 
\\ Ph.D. in Engeneering \hfill {\em April 2019 - March 2022} 
\\ Department of Advanced Interdisciplinary Studies
\\ Advisor: Dr. Masahiko Inami\\
\vspace{-3mm}
\\ {\bf Tohoku University} \hfill {\bf Sendai, Japan} 
\\ M.S. in Information Science \hfill {\em April 2017 - March 2019}
\\ Department of Applied information science
\\ Advisor: Dr. Masashi Konyo\\
\vspace{-3mm}
\\{\bf Tokyo University of Science} \hfill {\bf Tokyo, Japan} 
\\ B.S. in Mechanical Engineering\hfill {\em April 2013 - March 2017} 
\\ Department of Mechanical Engineering
\\ Advisor: Dr. Hiroshi Mizoguchi\\
%Minor in Linguistics \smallskip \\
%Member of Eta Kappa Nu \\
%Member of Upsilon Pi Epsilon \\


\end{rSection}

\begin{rSection}{Work Experience}
{\bf commissure Inc.} \hfill {\bf Tokyo, Japan} 
\\{\em CTO}\hfill {\em January 2023 - present} 

{\bf Keio University Graduate School of Media Design, Embodied Media} \hfill {\bf Tokyo, Japan} 
\\{\em Project Assistant Professor}\hfill {\em April 2023 - present} 

{\bf The University of Tokyo, Research Center for Advanced Science and Technology, Information Somatics Lab} \hfill {\bf Tokyo, Japan} 
\\{\em Project Assistant Professor}\hfill {\em April 2022 - March 2023} 
% \begin{itemize}
% \setlength{\itemsep}{-2mm}
%   \item Developing the full-body scale haptic device using rotational tactors array.
%   \item Investigating the perceptual effects by psychophysical experiments using developed devices. 
%   \item Creating the stimuli design method to generate the spatio-temporal haptic presentation.
% \end{itemize}

{\bf Information Somatics Lab, The University of Tokyo} \hfill {\bf Tokyo, Japan} 
\\{\em  Research Assistant / JSPS Research Fellow DC2, Advisor: Masahiko Inami}\hfill {\em April 2019 - March 2022} 
% \begin{itemize}
% \setlength{\itemsep}{-2mm}
%   \item Developing the full-body scale haptic device using rotational tactors array.
%   \item Investigating the perceptual effects by psychophysical experiments using developed devices. 
%   %\item Analyzing strain energy distribution of superelastic material using FEM analysis.
%   \item Creating the stimuli design method to generate the spatio-temporal haptic presentation.
% \end{itemize}

{\bf exiii } \hfill {\bf Tokyo, Japan} 
\\{\em Haptic Research Internsip}\hfill  {\em April 2019 - September 2019} 
% \begin{itemize}
% \setlength{\itemsep}{-2mm}
%   \item Developing the haptic presentation method for AR interaction by presenting force on a wrist  
% \end{itemize}

{\bf Human-Robot Informatics Lab, Tohoku University } \hfill {\bf Sendai, Japan} 
\\{\em  Research Assistant, Advisor: Masashi Konyo}\hfill {\em April 2017 - March 2019} 
% \begin{itemize}
% \setlength{\itemsep}{-2mm}
%   \item Developing the slip sensation presentation method using asymmetric vibration and high-frequency vibration together.
%   \item Developing the skin-stretching device for buttocks.
%   \item Investigating the relationship between hole-body self-motion perception and the force on buttocks
% \end{itemize}
\end{rSection}
%--------------------------------------------------------------------------------
%    Projects And Seminars
%-----------------------------------------------------------------------------------------------

\begin{rSection}{Grants}
% \\{\bf Collaborative business with Ishikawa prefecture, "Design of Self-other Inseparable Intercorporeality Through Body Surface Deformation Device", RCAST and ISCO} \hfill {\bf 2021 - 2023}
{\bf ACT-X Hardware in Future for Resilience of Real Space, "Design of Self-other Inseparable Intercorporeality Through Body Surface Deformation Device", JST} \hfill {\bf 2021 - 2023}
\\{\em Japan Science and Technologies} \\
\vspace{-3mm}

{\bf Japan Society for Promotion of Science DC2, "Self-motion Sensation Induction by the Presentation of Distributed Force", JSPS} \hfill {\bf 2020 - 2021}
\\{\em Japan Society for Promotion of Science} \\
\vspace{-3mm}
\end{rSection}

\begin{rSection}{Honors and Awards}
% \\{\bf ACT-X Hardware in Future for Resilience of Real Space, JST} \hfill {\bf 2021 - 2024}
% \\{\em Japan Science and Technologies} \\
% \vspace{-3mm}
{\bf Best Paper Award, CHI 2025}\hfill {\bf 2025}
\\{\em Co-author, "“It’s Like Being On Stage”: Conveying Dancers’ Expressiveness Through A Haptic-Installed Contemporary Dance Performance"} \\
\\{\bf Best Demo Award, SIGGRAPH Asia 2025 Emerging Technologies}\hfill {\bf 2025}
\\{\em Co-author, "Experience Sharing Box: A Platform for Capturing, Storing, and Playback of Multisensory Memory"} \\
\\{\bf Official Selection for Laval Vitual, SIGGRAPH 2024 Emerging Technologies}\hfill {\bf 2024}
\\{\em Co-author, "FEELTECH Wear: Enhancing Mixed Reality Experience with Wrist to Finger Haptic Attribution"} \\
\\{\bf Best Demo Award, EuroHaptics Conference 2024}\hfill {\bf 2024}
\\{\em Co-author, "Rotational Skin-stretch Distribution Creates Directional Force Sensation on the Wrist"} \\
\\{\bf Dean's Award, Graduate School of Engineering}\hfill {\bf 2022}
\\{\em The University of Tokyo} \\
% \vspace{-3mm}
\\{\bf Fellowships from Japan Society for Promotion of Science DC2, JSPS} \hfill {\bf 2020 - 2021}
\\{\em Japan Society for Promotion of Science} \\
\vspace{-3mm}
\\{\bf Best Paper Award, UbiComp MIMSVAI 2021}\hfill {\bf 2021}
\\{\em Co-author, "High-speed non-contact thermal display using infrared rays and shutter mechanism"} \\
\\{\bf Best Paper Award Nominated, World Haptics Conference 2021}\hfill {\bf 2021}
\\{\em Author, "Two-dimensional moving phantom sensation created by rotational skin stretch distribution"} \\
\vspace{-3mm}
\\{\bf Outstanding Performance Award, Asia Digital Art Award 2020} \hfill {\bf 2020}
\\{\em Author, "TorsionCrowds"} \\
\vspace{-3mm}
\\{\bf Disruptive Challenge Finalist, INNO-Vation 2019} \hfill {\bf 2019}
\\{\em Author, Ministry of Internal Affairs and Communications in Japan} \\
\vspace{-3mm}
\\{\bf Best Demonstration Award Gold Winner, AsiaHptics 2018} \hfill {\bf 2018}
\\{\em Author, "Enhancing Haptic Experience in a Seat with Two-DoF Buttock Skin Stretch"} \\
\vspace{-3mm}
\\{\bf Ultrahaptics Student Challenge Finalist, EuroHaptics 2018} \hfill {\bf 2018}
\\{\em Co-author, "Weather Teleport"} \\
\vspace{-3mm}
\\{\bf Robomech Award, Robotics and Mechatoronics 2018}\hfill {\bf 2018}
\\{\em Author, "Induction of acceleration sensation of self-motion by presenting shear force to buttocks"} \\
\vspace{-3mm}
\\{\bf Young Researcher's Award, VRSJ 2018}\hfill {\bf 2018}
\\{\em Author, "Presenting a pushing up feeling from the seat by buttocks skin stretch"} \\
\vspace{-3mm}
\\{\bf Excellent Presentation Award, SICE SI 2018}\hfill {\bf 2018}
\\{\em Author, "2-DoF Buttock Skin Stretch for Inducing Acceleration Perseption of Self-motion"} \\
\vspace{-3mm}
\\{\bf COLOPL Award, International collegiate Virtual Reality Contest 2017}\hfill {\bf 2017}
\\{\em Team Manager, "Shall we coffee cup?"} \\
\vspace{-3mm}
\\{\bf Young Researcher's Award, SICE SI 2017}\hfill {\bf 2017}
\\{\em Author, "Sliding direction and speed presentation Combining high frequency and asymmetric vibration"} \\
\vspace{-3mm}
\\{\bf Exellent Presentation Award, SICE SI 2017}\hfill {\bf 2017}
\\{\em Author, "Sliding direction and speed presentation
Combining high frequency and asymmetric vibration"} \\
\vspace{-3mm}


\end{rSection}
%----------------------------------------------------------------------------------------
%	TECHNICAL STRENGTHS SECTION
%----------------------------------------------------------------------------------------

\begin{rSection}{Peer-Reviewed Publications}
\begin{Subsection}{Journal Paper}
  \begin{enumerate}
    \item \bibentry{horie2020design}
    \item \bibentry{saito2021transparency}(Co-first author)
  \end{enumerate}
\end{Subsection}
\begin{Subsection}{Conference Paper}
  \begin{enumerate} 
  \item \bibentry{emiri2025piano}
  \item \bibentry{shen2025s}\textbf{(Best Paper Award)}
  \item \bibentry{umehara2024rotational}
  \item \bibentry{horie2023wearable}
  \item \bibentry{murata2023dynamic}
  \item \bibentry{ichihashi2021high}\textbf{(Best Paper Award)}
    \item \bibentry{horie2021Two}\textbf{(Nominated to Best Paper Award)}
    \item \bibentry{horie2021Encountered}
    \item \bibentry{shimobayashi2021}
    \item \bibentry{tanaka2020Dual}
    \item \bibentry{horie2018buttock}
  \end{enumerate}
\end{Subsection}
\begin{Subsection}{Demo}
  \begin{enumerate}  
    \item \bibentry{higashi2025experience}{ \bf SIGGRAPH Asia 2025 Emerging Technologies Best Demo Award}
    \item \bibentry{umehara2024feeltech}{ \bf SIGGRAPH 2024 Emerging Technologies Official Selection for Laval Vitual}
    \item \bibentry{tsujita2024haptoroom}
    \item \bibentry{umehara2024rotational}{ \bf Best Demo Award}
    \item \bibentry{taguchi2023multichannel} 
    \item \bibentry{uyama2023somatic}
    \item \bibentry{horie2022seeing}
    \item \bibentry{ichihashi2022themoblinds}
    \item \bibentry{horie2020Torsion}
    \item \bibentry{horie2019whole}
    \item \bibentry{horie2018enhancing}{ \bf Best Demo Award Gold Winner}
    \item \bibentry{horie2018buttock}
  \end{enumerate}
\end{Subsection}
\end{rSection}
\begin{rSection}{Non-Reviewed Publications (First Author)}
  \begin{enumerate}
    \item \bibentry{horie2019VRSJ}
    \item \bibentry{horie2018SI}{ \bf Excellent Presentation Award}
    \item \bibentry{horie2018VRSJ}{ \bf Young Researcher's Award}
    \item \bibentry{horie2018Robomech}{ \bf Robomech Award}
    \item \bibentry{horie2017SI}{ \bf Excellent Presentation Award,}{ \bf Young Researcher's Award}
    \item \bibentry{horie2017VRSJ}
  \end{enumerate}
\end{rSection}
%----------------------------------------------------------------------------------------
%	WORK EXPERIENCE SECTION
%----------------------------------------------------------------------------------------

\begin{rSection}{Patent}

{\bf 力覚提示装置及び力覚提示方法} \hfill {\bf 特願2019-163973} 
\\{\em 堀江新（75/100）, 稲見昌彦（25/100）} \hfill {\em }

{\bf 力覚提示装置及び力覚提示方法} \hfill {\bf PCT/JP2020/032185} 
\\{\em 堀江新（75/100）, 稲見昌彦（25/100）} \hfill {\em }

{\bf 運動教示装置、運動教示システム、運動教示方法およびプログラム} \hfill {\bf 特願2020-219284} 
\\{\em 堀江新（45/100）,檜山敦（35/100）, 稲見昌彦（20/100）} \hfill {\em }

\end{rSection}

\begin{rSection}{Teaching Experience}

{\bf Otherness and Interaction Technology in Embodied Interaction} 
\\{\bf Keio University Graduate School of Media Design} 
\\{\em Guest Lecturer} \hfill {\em November 2023 - Present}

{\bf Digital Interaction Prototyping} 
\\{\bf Keio University Graduate School of Media Design} 
\\{\em Lecturer} \hfill {\em May 2023 - Present}

{\bf Data Processing}
\\{\bf Tohoku University} 
\\{\em Teaching Assistant} \hfill {\em October 2017 - February 2018}

\end{rSection}

\end{document}
